\documentclass[12pt]{extarticle}

% ============================================================
% PACKAGES
% ============================================================

% Encodage et langue
\usepackage[french]{babel}

% Mise en page
\usepackage[a4paper, margin=2cm]{geometry}

\usepackage[francais,bloc,completemulti]{automultiplechoice}    

% Mathématiques
\usepackage{amsmath}
\usepackage{amssymb}

% Dessins
\usepackage{tikz}

% Listes
\usepackage{enumitem}

% Boîtes
\usepackage[most]{tcolorbox}

% Images
\usepackage{graphicx}

% Mise en page avancée
\usepackage{fancyhdr}
\usepackage{multicol}

% Tableaux
\usepackage{booktabs}
\usepackage{tabularx}
\usepackage{colortbl}

\usepackage[normalem]{ulem}

\usepackage{pgfplots}
\pgfplotsset{compat=1.18}
\usepackage{pgf-pie}


\usepackage{ifthen}

% ============================================================
% COULEURS
% ============================================================

% ------------------------------------------------------------
% Couleurs principales
% ------------------------------------------------------------

\definecolor{Bleu}{RGB}{40, 100, 180}
\definecolor{Vert}{RGB}{50, 150, 90}
\definecolor{Orange}{RGB}{230, 140, 40}

\definecolor{BleuClair}{RGB}{235, 245, 255}
\definecolor{VertClair}{RGB}{235, 250, 240}
\definecolor{OrangeClair}{RGB}{255, 245, 225}

% ------------------------------------------------------------
% Couleurs complémentaires
% ------------------------------------------------------------

% Rappels de cours — violet
\definecolor{Violet}{RGB}{125, 85, 170}
\definecolor{VioletClair}{RGB}{245, 238, 252}

% Théorèmes — rouge
\definecolor{Rouge}{RGB}{190, 55, 55}
\definecolor{RougeClair}{RGB}{255, 238, 238}

% Lemme — indigo
\definecolor{Indigo}{RGB}{75, 85, 155}
\definecolor{IndigoClair}{RGB}{240, 241, 252}

% Points méthode — turquoise
\definecolor{Turquoise}{RGB}{35, 145, 150}
\definecolor{TurquoiseClair}{RGB}{235, 250, 250}

% Erreurs / attention — jaune
\definecolor{Jaune}{RGB}{210, 165, 35}
\definecolor{JauneClair}{RGB}{255, 250, 225}

% ------------------------------------------------------------
% Neutres
% ------------------------------------------------------------

\definecolor{GrisClair}{RGB}{245, 245, 245}
\definecolor{Gris}{RGB}{100, 100, 100}


% ============================================================
% COMMANDES PERSONNALISÉES
% ============================================================

\newcommand{\bb}[1]{\textbf{#1}}
\newcommand{\ii}[1]{\textit{#1}}

\newcommand{\cm}{\,\mathrm{cm}}
\newcommand{\m}{\,\mathrm{m}}

\newcommand{\reponse}[1]{\boxed{#1}}

% ============================================================
% TABLEAUX
% ============================================================

\newcommand{\ligneentete}{%
    \rowcolor{Bleu}%
    \color{white}\bfseries
}

\newcommand{\entete}[1]{%
    \cellcolor{Bleu}\color{white}\bb{#1}%
}


\NewDocumentCommand{\hl}{O{JauneClair}m}{%
  \tcbox[
    on line,
    colback=#1,
    colframe=#1,
    boxrule=0pt,
    arc=4pt,
    left=0pt,
    right=0pt,
    top=1pt,
    bottom=1pt
  ]{#2}%
}

\NewDocumentCommand{\hlb}{O{JauneClair}m}{%
  \tcbox[
    on line,
    colback=#1,
    colframe=#1,
    boxrule=0pt,
    arc=4pt,
    left=0pt,
    right=0pt,
    top=1pt,
    bottom=1pt
    ]{\bb{#2}}%
}
\newlength{\answerwidth}


\newcommand{\answer}[2][VertClair]{%
    \ifmmode
        \settowidth{\answerwidth}{$\displaystyle #2$}%
        \ifanswers
            {\setlength{\fboxsep}{0pt}%
             \colorbox{#1}{$\displaystyle #2$}}%
        \else
            \makebox[0pt][l]{%
                \makebox[\answerwidth][l]{%
                    \leaders\hbox to 0.5em{\hss.\hss}\hfill
                }%
            }%
            \phantom{#2}%
        \fi
    \else
        \settowidth{\answerwidth}{#2}%
        \ifanswers
            \hl[#1]{#2}%
        \else
            \makebox[\answerwidth][l]{%
                \leaders\hbox to 0.5em{\hss.\hss}\hfill
            }%
        \fi
    \fi
}


\newcommand{\graphanswer}[1]{%
    \ifanswers
        #1%
    \else
        0%
    \fi
}

\newcommand{\exo}[1]{%
    \includegraphics[
        width=\linewidth,
        keepaspectratio
    ]{#1}%
}

\newcommand{\nodata}{}
\newcommand{\correctanswer}[1]{
    \explain{
            \newline\vspace{10pt}
            #1
    }
}

\newlength{\mylen}
\newtcolorbox{exercicebox}[1][\unskip]{
    colback=white,
    colframe=Bleu,
    boxrule=1pt,
    arc=3mm,
    title=\bb{Exercice #1},
    left=2mm,
    right=2mm,
    top=2mm,
    bottom=1mm
}

\newcounter{theorem}
\newtcolorbox{theoremebox}[1][\unskip]{
    colback=RougeClair,
    colframe=Rouge,
    boxrule=1pt,
    arc=3mm,
    title=\bb{Théorème \stepcounter{theorem}\thetheorem}\hspace{5pt}#1,
    left=5mm,
    right=5mm,
    top=4mm,
    bottom=4mm
}

\newcounter{probleme}
\newtcolorbox{probbox}[1][\unskip]{
    colback=white,
    colframe=Rouge,
    boxrule=1pt,
    arc=3mm,
    title=\bb{Problème #1},
    left=5mm,
    right=5mm,
    top=2mm,
    bottom=2mm
}

\newenvironment{tableau}[1]
{%
    \renewcommand{\arraystretch}{1.3}%
    \begin{center}%
    \begin{tabularx}{\linewidth}{#1}%
    \toprule
}
{%
    \bottomrule
    \end{tabularx}%
    \end{center}%
    \renewcommand{\arraystretch}{1}%
}



% Set page layout after the config files
\geometry{
  paper=a4paper,
  hmargin=1cm,
  headheight=1cm,
  headsep=.3cm,
  footskip=0cm,
  top=1.5cm,
  left=1.3cm,
  right=1.3cm,
  bottom=1.5cm
}

%\geometry{hmargin=1cm,headheight=1cm,headsep=.3cm,footskip=0cm,top=1.5cm,bottom=1.5cm}

\def\AMCotextReserved{\emph{(Reservé)}}
\AMCopenOpts{boxframerule=0pt,boxframerulecol=red,framerule=0pt}
\newdimen\AMCpostOquest\AMCpostOquest=-7mm
\AMCformHSpace=.3em
\AMCinterBquest=1ex

\makeatletter
\def\jepasse{%
\def\AMCbeginQuestion##1##2{\ifAMC@catalog\textbf{##1}\fi}
\AMCquestionNumberfalse}
\makeatother


\begin{document}

\exemplaire{24}{    

% ---- Header ------
\fbox{
    \begin{minipage}{.40\linewidth}
        \namefield{Nom: }
        Prénom :\\
        Date:
    \end{minipage}
}
\hfill
\begin{minipage}{.48\linewidth}
    \AMCboxStyle{size=1.7ex,down=.2ex}
    \AMCcodeGridInt[vertical=false]{etu}{2}
\end{minipage}
% ---- Fin Header ------


\newcommand{\hpt}[1]{
    \AMCboxStyle{shape=square,color=black,size=1.8ex,down=.2ex}
    \begin{minipage}{2ex}
        \jepasse
        \begin{question}{##1}\bareme{b=1}
            \begin{reponseshoriz}
                \bonne{}
            \end{reponseshoriz}
        \end{question}
    \end{minipage} 
}
\newcommand{\pt}[1]{
    \AMCboxStyle{shape=oval,color=black,size=2ex,down=.2ex}
    \begin{minipage}{2ex}
        \jepasse
        \begin{question}{##1}\bareme{b=1}
            \begin{reponseshoriz}
                \bonne{}
            \end{reponseshoriz}
        \end{question}
    \end{minipage} 
}
 
 \begin{exercicebox}[\bb{1}\hspace{3mm}\it{(Questions de cours)}
    \hfill\pt{cours1}\pt{cours2}\pt{cours3}\pt{cours4}\hpt{cours4.5}\hpt{cours5}]
    \begin{enumerate}
        \item Donne la définition du carré d'un nombre.
        \item Donne la définition de la racine carrée d'un nombre.
        \item Enonce le théorème de Pythagore.
        \item Donne la racine carrée des nombres suivants : 64, 36, 121.
        \item Un élève affirme qu'une racine carré de 16 est -4 car
            $(-4)\times(-4)=16$. A-t-il raison ? Justifie.
    \end{enumerate}
\end{exercicebox}

 \begin{exercicebox}[\bb{2}\hspace{3mm}\it{(Calculer des longueurs)}
    \hfill\pt{op1}\pt{op2}\pt{op3}\pt{op4}\pt{op5}\pt{op6}\pt{op7}\pt{op8}\pt{op9}\hpt{op9.5}\hpt{op10}]
    \begin{enumerate}
        \item Pour chacun des triangles suivants, calcule la longueur manquante,
            d'abord en valeur exacte, puis en valeur approchée au cm.
            \bb{Détaille ton raisonnement}.
        \begin{enumerate}
            \item Le triangle ABC est rectangle en A, AB=9m et AC=8m. Calcule BC.
            \item Le triangle DEF est rectangle en E, EF=3m et DF=11m. Calcule ED.
        \end{enumerate}
        \begin{minipage}{0.40\textwidth}
            \includegraphics[width=\textwidth]{images/angle.png}
        \end{minipage}
        \hfill
        \begin{minipage}{0.49\textwidth}
            \item Dans le triangle ci-contre
                \begin{enumerate}
                    \item Justifie que le triangle est rectangle en C.
                    \item Calcule la longueur BC, en valeur exacte puis arrondie au
                        millimètre.
                \end{enumerate}
        \end{minipage}
    \end{enumerate}
\end{exercicebox}


\begin{probbox}[\bb{1}
     \hfill\pt{prob1}\pt{prob2}\pt{prob3}\pt{prob4}\hpt{prob4.5}\hpt{prob5}
    ]
    \begin{minipage}{0.58\textwidth}
        Pour apprendre son métier, un apprenti
        maçon a monté un mur en briques, de 0,90 m de
        hauteur. Son patron arrive pour vérifier son
        travail : il marque un point B sur le mur, à 80 cm
        du sol, et un point A, à 60 cm du pied du mur. Il
        mesure alors la distance entre les points A et B et
        obtient 1 m.
        \begin{enumerate}
            \item L'apprenti a-t-il construit un mur parfaitement
        perpendiculaire au sol ? Justifie.
        \end{enumerate}


    \end{minipage}
    \hfill
    \begin{minipage}{0.38\textwidth}
        \includegraphics[width=\textwidth]{images/mur.png}
    \end{minipage}
\end{probbox}





%\begin{probbox}
%    Un groupe d'amis décident de faire une excursion dont le prix est fixé par
%    une agence de voyage.
%
%    \begin{itemize}
%        \item 1er cas : si chacun des amis verse 150€, il manque alors 1500€ pour
%            l'excursion.
%        \item 2eme cas : si chacun des amis vers 200€, il y a 200€ de trop.
%    \end{itemize}
%
%    Quel est le nombre de participant à l'excursion ?\\
%    Quel est le prix fixé par l'agence de voyage pour l'excursion ?
%\end{probbox}

}   
\end{document}
